\section*{Overview}

Suffix automaton is the suffix structure I want when the string is built online or when substring-state transitions feel
more natural than suffix-array order. It represents all substrings of one string in a DAG with only \(O(n)\) states.

\section*{When to Use It}

Use a suffix automaton when:

\begin{itemize}
  \item distinct substrings, occurrence counts, or longest common substring style queries appear,
  \item the string is processed left to right,
  \item an automaton over substring states feels easier than sorting suffixes.
\end{itemize}

\section*{Core Idea}

Each state represents an equivalence class of substrings that share the same set of ending positions. The automaton is
extended one character at a time, and occasionally a clone state is created to preserve the correct suffix-link and
transition structure.

\section*{Key Insight}

The suffix link points to the largest proper suffix class that still matters after extending the string. That is why the
automaton remains linear even though it compactly represents every substring.

\section*{Operations / Main Technique}

\begin{itemize}
  \item extend by one character,
  \item follow transitions to test whether a pattern is a substring,
  \item use \(\texttt{len[state] - len[link[state]]}\) to count distinct substrings,
  \item propagate end-position counts in topological length order for occurrence queries.
\end{itemize}

\paragraph{Worked example.}
Every time a new character is appended, all new substrings are exactly the suffixes ending at that new position. The
automaton extension adds those efficiently without rebuilding older states.

\section*{Correctness Intuition}

The construction maintains the smallest automaton that recognizes all suffixes of the processed prefix, which is enough
to recognize all substrings as well. Clone states appear only when an old state needs to be split into two equivalence
classes with different longest-length behavior.

\section*{Complexity Analysis}

The suffix automaton is built in \(O(n)\) time and memory for a fixed alphabet representation that supports \(O(1)\)
transitions on average or with arrays.

\section*{Implementation}

The code supports:

\begin{itemize}
  \item online extension,
  \item substring existence checks,
  \item counting distinct substrings,
  \item occurrence counts after one propagation pass.
\end{itemize}

\section*{Common Pitfalls}

\begin{itemize}
  \item Getting clone-state initialization wrong.
  \item Forgetting that occurrence counts need a later propagation step.
  \item Using a suffix automaton when the problem is really about lexicographic order, where suffix array fits better.
\end{itemize}

\section*{Variants / Extensions}

\begin{itemize}
  \item Longest common substring between strings.
  \item Build the suffix-link tree for further DP.
  \item Palindromic tree when the structure is about palindromes rather than arbitrary substrings.
\end{itemize}

\section*{Practice Problems}

\begin{itemize}
  \item Count distinct substrings.
  \item Query whether many patterns are substrings.
  \item Longest common substring or repeated-substring tasks.
\end{itemize}

\section*{References}

\begin{itemize}
  \item \href{https://cp-algorithms.com/string/suffix-automaton.html}{cp-algorithms: Suffix Automaton}.
  \item \href{https://usaco.guide/adv/string-suffix}{USACO Guide: String Suffix Structures}.
  \item \href{https://codeforces.com/catalog?locale=en}{Codeforces Catalog}.
\end{itemize}
